%% ============================================================================
\section{Conclusion}
\label{sec:conclusion}
%% ============================================================================

This paper has presented a controlled study in which wideband acoustic emission
and heterodyned passive ultrasound were acquired simultaneously from a ball
bearing across staircase speed sweeps and stepped temperature blocks, and
regressed to the ISO~281 viscosity ratio $\kappa$ under an acquisition-grouped,
leakage-controlled evaluation protocol. The main findings are as follows:

\begin{enumerate}
  \item \textbf{The $\kappa$ regression acts as an operating-point soft sensor.}
  A two-stage decomposition that predicts speed and temperature from the features and then calculates $\kappa$ recovers essentially all of
  the direct regression performance for every channel, leaving no measurable lubrication-specific residual. Reported accuracies for acoustic $\kappa$ regression, here and in prior work, should be interpreted as operating-point recovery. Detecting lubrication changes that are not caused by the operating point through acoustic measurements remains an open problem.
  \item \textbf{AE and passive US are complementary.} The combined feature set outperforms either channel alone for both model families,providing quantitative indication of the complementarity. The complementarity has a structural basis: both channels track rotational speed, but US resolves temperature substantially better than AE. Which single channel performs best is model-family dependent, so single-model channel comparisons should be interpreted with caution.
  \item \textbf{Leakage-controlled evaluation changes the picture.}
  Grouping segments by acquisition hold and operating-point similarity prevents near-duplicate windows from crossing the training--evaluation boundary, and group-level significance testing avoids the inflation of window-level tests. The reported figures are conservative, deployment-representative estimates.
\end{enumerate}
